% !PRINTED 523
\weberpage{523}

%<<CONT
\gesperrt{cienten ganze Zahlen in $\Omega$ sind, und geht durch Multiplication mit einem rationalen Functional, welches eine Einheit ist, in $\chi$ \"uber.}

\begin{center}
\S.\,143.

\gesperrt{Die Primfactoren der Zahlen des K\"orpers $\Omega$.}
\end{center}

Da unter den Functionalen des K\"orpers $\Omega$ auch die Zahlen enthalten sind, so ergiebt sich aus den Resultaten des \S.\,141, dass sich auch die ganzen Zahlen des K\"orpers $\Omega$ in Primfactoren zerlegen lassen, aber in Primfactoren, die im Allgemeinen nicht Zahlen, sondern Functionale sind. Es ist aber nicht ausgeschlossen, dass in besonderen F\"allen unter den Primfunctionalen auch Zahlen auftreten k\"onnen, die dann \gesperrt{Primzahlen des K\"orpers $\Omega$} heissen.

Alle Gleichungen zwischen Functionalen sind in letzter Instanz identische Gleichungen zwischen ganzen rationalen Functionen und l\"osen sich in eine Reihe von Gleichungen zwischen Zahlen auf. Sie bleiben also richtig, wenn f\"ur die Variablen andere Zeichen gesetzt werden.

Es folgt daraus, dass ein ganzes Functional, eine Einheit, ein Primfunctional nicht aufh\"oren, ganze Functionale, Einheiten, Primfunctionale zu sein, wenn f\"ur die Variablen andere Symbole gesetzt werden.

Zerlegen wir also eine Zahl in ihre Primfactoren, so k\"onnen, wenn unter diesen Primfactoren Functionale vorkommen, in der diese Zerlegung darstellenden Gleichung die Variablen durch beliebige andere Variable ersetzt werden, und daraus folgt die Verallgemeinerung:

1. \gesperrt{Man kann die Primfactoren eines ganzen Functionals $\omega$ so darstellen, dass sie die Variablen, von denen $\omega$ abh\"angt, nicht enthalten.}

Denn jedes ganze Functional $\omega$ ist Theiler von Zahlen, sogar von rationalen Zahlen, z. B. von der absoluten Norm von $\omega$. Die Primfactoren von $\omega$ sind daher unter den Primfactoren einer dieser Zahlen zu suchen und k\"onnen also durch Variable dargestellt werden, die von den in $\omega$ vorkommenden verschieden sind.
